\section{Related Work}
\label{sec:related}

The relevant literature spans personalized language generation, user modeling,
implicit feedback, adaptive interfaces, recommender systems, conversational
memory, persona conditioning, and preference optimization. Context-aware
recommender systems and implicit-feedback methods established that behavior is
not equivalent to positive preference and that context belongs in the utility
model \citep{hu-etal-2008-collaborative,adomavicius-2011-context-aware}.

Recent LLM personalization work has moved beyond fixed persona sentences. LaMP
established retrieval-based personalization tasks
\citep{salemi-etal-2024-lamp}; other systems use per-user adapters or embeddings
\citep{tan-etal-2024-democratizing,liu-etal-2025-llms}, domain-aware compressed
profiles \citep{su-etal-2025-personalized}, and latent preference attributes
\citep{li-etal-2025-prefpalette}. PLUS jointly learns an interpretable natural-
language summary and a personalized reward model \citep{nam-etal-2026-plus}, while
S2Pref distinguishes stable from situational preferences and includes preference
conflicts \citep{gao-etal-2026-beyond}. These results mean that dynamic profiles,
domain conditioning, and interpretable summaries are not individually novel.

Evaluation work also warns against assuming that personalization helps. PrefEval
tests explicit and implicit preference following in long contexts
\citep{zhao-etal-2025-prefeval}, and PrefDisco reports settings where uninformed
personalization underperforms a generic response \citep{li-etal-2026-prefdisco}.
Broad memory systems offer temporal retrieval and compression mechanisms
\citep{park-etal-2023-generative,packer-etal-2023-memgpt,chhikara-etal-2025-mem0}
but usually optimize episodic or factual recall rather than an editable response-
preference policy.

The narrower gap investigated here is the joint system of typed response
dimensions, hierarchical applicability, calibrated evidence provenance, user
governance, provider-neutral compilation, and behavioral cross-model evaluation.
The structured literature review found no peer-reviewed demonstration of the
entire bundle at its evidence cutoff; this is an absence-of-evidence statement,
not a priority claim. The review protocol and study selection limitations are
reported in the accompanying artifact.
